\documentclass[10pt]{article}
\usepackage[T1]{fontenc}
\usepackage[utf8]{inputenc}
\usepackage[french]{babel}
\usepackage{amsmath,amssymb,mathrsfs,mathpazo}
\usepackage{geometry}
\geometry{a4paper,margin=2cm}
\usepackage{graphicx,array,tabularx,tabularray,colortbl,multirow,multicol,booktabs,enumitem,pifont,ulem,fancybox,diagbox,cancel}
\usepackage{pstricks,pstricks-add,pst-plot,pst-tree,pst-node,pst-text,pst-func,pst-math,pst-bspline,pst-3dplot,pst-eucl,pst-solides3d}
\let\pscalculate\undefined
\usepackage{pst-bezier}
\usepackage{tikz,pgfplots,tkz-tab}
\usetikzlibrary{arrows,arrows.meta,patterns,calc,positioning,trees,decorations.markings,shapes.geometric,shapes.misc,shapes.arrows,intersections,decorations.pathmorphing,decorations.pathreplacing,backgrounds,shadows,babel}
\pgfplotsset{compat=1.18}
\usepgfplotslibrary{fillbetween}
\usepackage[np]{numprint}
\usepackage{listings,makecell,pas-tableur,twoopt}
\let\barre\undefined
\usepackage[pstricks]{tablvar}
\definecolor{darkgreen}{RGB}{0,100,0}
\providecommand{\red}{\color{red}}
\providecommand{\blue}{\color{blue}}
\providecommand{\green}{\color{green}}
\providecommand{\black}{\color{black}}
\providecommand{\white}{\color{white}}
\providecommand{\gray}{\color{gray}}
\providecommand{\cyan}{\color{cyan}}
\providecommand{\magenta}{\color{magenta}}
\providecommand{\yellow}{\color{yellow}}
\providecommand{\violet}{\color{violet}}
\providecommand{\orange}{\color{orange}}
\providecommand{\euro}{\texteuro}
\providecommand{\EUR}{\texteuro}
\providecommand{\ieme}{\textsuperscript{e}}
\providecommand{\degres}{\ensuremath{^\circ}}
\providecommand{\degree}{\ensuremath{^\circ}}
\providecommand{\vv}[1]{\overrightarrow{#1}}
\providecommand{\Coord}[1]{#1}
\providecommand{\decosix}{}
\providecommand{\starredbullet}{\ensuremath{\bullet}}
\pagestyle{empty}
\setlength{\parindent}{0pt}
\usepackage{adjustbox}\DeclareUnicodeCharacter{2236}{:}\DeclareUnicodeCharacter{2190}{\ensuremath{\leftarrow}}
\providecommand{\R}{\mathbb{R}}
\providecommand{\N}{\mathbb{N}}
\providecommand{\Z}{\mathbb{Z}}
\providecommand{\Q}{\mathbb{Q}}
\providecommand{\D}{\mathbb{D}}
\providecommand{\C}{\mathbb{C}}
\providecommand{\e}{{\mathrm{e}}}
\providecommand{\vect}[1]{\overrightarrow{#1}}
\providecommand{\vectt}[1]{\overrightarrow{\mathrm{#1}}}
\providecommand{\vvt}[1]{\overrightarrow{\mathrm{#1}}}
\providecommand{\np}[2][]{{#2~#1}}
\providecommand{\vv}[1]{{\overrightarrow{#1}}}
\renewcommand{\d}{{\mathrm{d}}}
\providecommand{\degres}{{^\circ}}
\providecommand{\degre}{{^\circ}}
\providecommand{\blue}{\color{blue}}
\providecommand{\red}{\color{red}}
\providecommand{\enskip}{\quad}
\providecommand{\llbracket}{\lbrack\!\lbrack}
\providecommand{\rrbracket}{\rbrack\!\rbrack}
\providecommand{\euro}{{\text{€}}}
\providecommand{\no}{{\text{n°}}}
\providecommand{\interval}[2]{{[#1;#2]}}
\providecommand{\intervFO}[2]{{[#1;#2[}}
\providecommand{\dotfill}{\ldots}

\providecommand{\ds}{\displaystyle}
\providecommand{\textding}[1]{#1}
\providecommand{\Oij}{\(\mathrm{O};\vec{\imath},\vec{\jmath}\)}
\providecommand{\Oijk}{\(\mathrm{O};\vec{\imath},\vec{\jmath},\vec{k}\)}
\providecommand{\Ouv}{\(\mathrm{O};\vec{u},\vec{v}\)}
\providecommand{\intervFF}[2]{\ensuremath{[#1;#2]}}
\renewcommand{\euro}{\ifmmode\text{\texteuro}\else\texteuro\fi}
\DeclareUnicodeCharacter{00B0}{\ifmmode^{\circ}\else\textdegree\fi}
\DeclareUnicodeCharacter{20AC}{\ifmmode\text{\texteuro}\else\texteuro\fi}
% Énoncés : Éducation nationale. Transcription et corrigés : APMEP, Denis Vergès et contributeurs. Découpage et classement : CoopMaths, Sébastien Lozano, Éric Elter et contributeurs. Adaptation de la présentation : Maths Explorer.
% CC BY-SA : https://coopmaths.fr/www/about/
% https://coopmaths.fr/alea/static/bac/2021/tex/bac_2021_05_ameriquenord_4_cor.tex
\begin{document}
\noindent\textit{APMEP / CoopMaths. CC BY-SA. Présentation adaptée par Maths Explorer.}\par\medskip
\smallskip













\medskip

\textbf{Affirmation 1 :} Pour tous réels $a$ et $b$,\, $\left(\text{e}^{a+b}\right)^2 =
 \text{e}^{2(a + b)} =  \text{e}^{2a + 2b} = \text{e}^{2a} \times \text{e}^{2b}$. Donc l'affirmation est fausse.

\smallskip
\textbf{Affirmation 2 :}  
Une équation de la tangente $t$ au point A d'abscisse 0 à la courbe représentative de la fonction $f$ définie sur $\R$ par $f(x) = - 2 + (3 - x)\text{e}^x$ est :

$M(x~;~y) \in t \iff y - f(0) = f'(0)(x - 0)$.

Or $f(0) = - 2 + 3\text{e}^0 = - 2 + 3 = 1$ et 

$f'(x) = -\text{e}^x + (3 - x)\text{e}^x = (2 - x)\text{e}^x$, d'où $f'(0) = 2\text{e}^0 = 2$.

Donc $M(x~;~y) \in t \iff y - 1 = 2(x - 0) \iff y = 2x + 1$. L'affirmation est vraie.

\smallskip

\textbf{Affirmation 3 :} 
On a quel que soit le réel $x$ : $\text{e}^{2x} - \text{e}^{x} + \dfrac{3}{x}= \text{e}^{x}\left(\text{e}^{x} - 1 + \dfrac{3}{x\text{e}^{x}}\right)$.

On a $\displaystyle\lim_{x \to + \infty} x\text{e}^{x} = + \infty$, donc $\displaystyle\lim_{x \to + \infty} \dfrac{3}{x\text{e}^{x}} = 0$.

Par somme de limites $\displaystyle\lim_{x \to + \infty}\text{e}^{x} - 1 + \dfrac{3}{x\text{e}^{x}} = + \infty$ et par produit de limites :

$\displaystyle\lim_{x \to + \infty}\text{e}^{x}\left(\text{e}^{x} - 1 + \dfrac{3}{x\text{e}^{x}}\right) = + \infty$. L'affirmation est fausse.

\smallskip

\textbf{Affirmation 4 :} 
Soit la fonction $f$ définie sur $\R$ par $f(x) = 1 - x + \text{e}^{-x}$

$f$ somme de fonctions dérivables sur $\R$ est dérivable et sur cet intervalle :

$f'(x) = - 1 - \text{e}^{-x} = - \left(1 + \text{e}^{-x} \right) < 0$ car quel que soit le réel $x$, \, $\text{e}^{-x}  > 0$, donc $1 + \text{e}^{-x} > 1$ puis $- \left(1 + \text{e}^{-x} \right) < - 1 < 0$.

La fonction $f$ est strictement décroissante sur $\R$.



On a $\displaystyle\lim_{x \to - \infty} \text{e}^{-x} = + \infty$, et $\displaystyle\lim_{x \to - \infty} - x = + \infty$\, donc par somme de limites  $\displaystyle\lim_{x \to - \infty} f(x) = + \infty$.

De même $\displaystyle\lim_{x \to + \infty} - x = - \infty$ et $\displaystyle\lim_{x \to + \infty} \text{e}^{- x} = 0$, d'où par somme de limites : $\displaystyle\lim_{x \to + \infty} f(x) = - \infty$.

D'après le théorème des valeurs intermédiaires il existe donc $x_0 \in \R$ tel que $f\left(x_0\right) = 0$.

Comme $f(0) = 1 + 1 = 2 > 0$ et $f(2) = 1 - 2 + \text{e}^{-2} \approx - 0,86 < 0$, on a bien $0 < x_0 < 2$. L'affirmation est vraie.
\smallskip

\textbf{Affirmation 5 :} 
On a $g'(x) = 2x - 5 + \text{e}^x$ et $g''(x) = 2 + \text{e}^x > 0$ comme somme de deux termes supérieurs à zéro. la fonction $g$ est donc convexe sur $\R$. L'affirmation est vraie.

\bigskip

\end{document}
